% Converted from README.md
\section{Debagovanje}\label{debagovanje}

Debager se koristi kada treba posmatrati konkretno izvršavanje programa: vrednosti promenljivih, stek poziva, registre, niti, memoriju i tačku na kojoj se ponašanje odvaja od očekivanog. Pre debagovanja treba imati reprodukciju problema i, ako je moguće, izvršivu verziju prevedenu sa simbolima za debagovanje.

\subsection{Okruženje izvršavanja i strategija prevođenja}\label{okruzenje-izvrsavanja-strategija-prevodenja}

Pre izbora debagerskog pristupa treba znati kako program nastaje i gde se izvršava. Interpreter, JIT (Just-In-Time) i AOT (Ahead-Of-Time) okruženja različito mapiraju izvorni kod, međukod i mašinski kod. To utiče na breakpoint-e, simbole, stack trace, profajliranje, instrumentaciju i reprodukciju grešaka.

\definecolor{RuntimeGoodText}{HTML}{4F8F63}
\definecolor{RuntimeWarnText}{HTML}{B08A2B}
\definecolor{RuntimeBadText}{HTML}{B45F5F}
\newcommand{\rtgood}[1]{\textcolor{RuntimeGoodText}{#1}}
\newcommand{\rtwarn}[1]{\textcolor{RuntimeWarnText}{#1}}
\newcommand{\rtbad}[1]{\textcolor{RuntimeBadText}{#1}}

\noindent\rtgood{Zeleno} označava povoljan aspekt, \rtwarn{žuto} kompromis ili oprez, a \rtbad{crveno} nepovoljan aspekt za debagovanje, profajliranje ili isporuku.

\begingroup
\renewcommand{\arraystretch}{1.18}
\begin{longtable}{@{}p{0.30\linewidth}p{0.32\linewidth}p{0.30\linewidth}@{}}
\toprule\noalign{}
Interpreter & JIT & AOT \\
\midrule\noalign{}
\endhead
\bottomrule\noalign{}
\endlastfoot
\rtwarn{Tumači kod tokom izvršavanja.} & \rtwarn{Tumači, zatim prevodi vruće putanje.} & \rtwarn{Prevodi pre pokretanja.} \\
\addlinespace[0.25em]
\rtgood{Brz razvojni ciklus.} & \rtgood{Brz razvojni ciklus.} & \rtbad{Sporiji ciklus zbog punog prevođenja.} \\
\addlinespace[0.25em]
\rtwarn{Izvršava samo posećene putanje.} & \rtwarn{Kreće od bajtkoda ili međukoda.} & \rtwarn{Optimizuje ceo poznati program.} \\
\addlinespace[0.25em]
\rtgood{Greške se vide u runtime-u.} & \rtwarn{Runtime menja kod tokom rada.} & \rtbad{Optimizacije mogu zakloniti uzrok.} \\
\addlinespace[0.25em]
\rtgood{Prenosiv uz isti interpreter.} & \rtgood{Prenosiv uz isti runtime.} & \rtbad{Vezan za ciljnu platformu.} \\
\addlinespace[0.25em]
\rtbad{Sporije pokretanje.} & \rtbad{Sporije pokretanje i zagrevanje.} & \rtgood{Brzo pokretanje.} \\
\addlinespace[0.25em]
\rtbad{Niže vršne performanse.} & \rtgood{Visoke posle zagrevanja.} & \rtgood{Visoke uz dobre optimizacije.} \\
\addlinespace[0.25em]
\rtwarn{Profajl meri i interpreter.} & \rtgood{Runtime često ima dobru podršku.} & \rtbad{Teže zbog inlining-a i simbola.} \\
\addlinespace[0.25em]
\rtgood{Laka instrumentacija.} & \rtgood{Runtime hook-ovi pomažu.} & \rtbad{Instrumentacija pre prevođenja ili binarno.} \\
\addlinespace[0.25em]
\rtgood{Dinamičke osobine su prirodne.} & \rtgood{Dinamičke osobine su prirodne.} & \rtbad{Dinamika je obično ograničena.} \\
\addlinespace[0.25em]
\rtbad{Veći memorijski otisak.} & \rtbad{Runtime, JIT kod i metapodaci.} & \rtgood{Manji memorijski otisak.} \\
\addlinespace[0.25em]
\rtbad{Potreban interpreter/runtime.} & \rtbad{Potreban runtime.} & \rtgood{Runtime nije nužan kod statičkog povezivanja.} \\
\end{longtable}
\endgroup

\subsection{Kada koristiti debager}\label{kada-koristiti-debager}

\begin{longtable}{@{}p{0.26\linewidth}p{0.34\linewidth}p{0.32\linewidth}@{}}
\toprule\noalign{}
Situacija & Šta debager daje & Kada izabrati drugi alat \\
\midrule\noalign{}
\endhead
\bottomrule\noalign{}
\endlastfoot
Pad programa koji se ponavlja & Stek poziva, argumente funkcija i stanje memorije u trenutku pada. & Ako se pad ne ponavlja lokalno, prvo prikupiti core dump ili logove. \\
Pogrešna vrednost u toku izvršavanja & Watchpoint, breakpoint i korak-po-korak praćenje mesta gde se vrednost menja. & Ako postoji sumnja na prekoračenje memorije, prvo uključiti ASan ili Memcheck. \\
Zaglavljen proces & Stekove svih niti i lokaciju čekanja, na primer muteks, \texttt{futex}, I/O ili beskonačnu petlju. & Ako treba profil vremena čekanja kroz duži period, koristiti \texttt{perf}, eBPF ili runtime profajler. \\
Postmortem analiza & Analizu core dump-a bez ponovnog pokretanja programa. & Ako nema simbola i tačne binarne verzije, prvo obezbediti artefakte za reprodukciju. \\
Interakcija sa hardverom ili mapiranom memorijom & Registre, memorijske adrese, asemblerski prikaz i uticaj optimizacija. & Ako se problem vidi samo kao sistemski događaj, prvo koristiti \texttt{strace}, \texttt{dmesg} ili eBPF. \\
\end{longtable}

Tipičan tok rada je: pokrenuti program sa istim ulazom, zaustaviti se pre sumnjivog mesta, proveriti pretpostavke o stanju, izvršiti sledećih nekoliko koraka i uporediti očekivano stanje sa stvarnim. Debager ne dokazuje da greška ne postoji; on objašnjava jedno konkretno izvršavanje.

\subimport{01_gdb/}{index.tex}

\subimport{02_java/}{index.tex}
